\documentclass[12pt,a4paper]{article}
 
% Pacotes essenciais de formatação (mesmos do TP1) 
\usepackage[utf8]{inputenc}
\DeclareUnicodeCharacter{201C}{``}
\DeclareUnicodeCharacter{201D}{''}
\usepackage[T1]{fontenc}
\usepackage[brazil]{babel}
\usepackage{mathptmx}
\usepackage[left=3.0cm, top=3.0cm, right=2.0cm, bottom=2.0cm]{geometry}
\usepackage{setspace}
\onehalfspacing
 
\usepackage{indentfirst}
\setlength{\parindent}{1.25cm}
\setlength{\parskip}{0pt}
 
\usepackage{graphicx}
\graphicspath{{prints/}}
\usepackage{xcolor}
\usepackage{amsmath,amssymb}
\usepackage{listings}
\usepackage{fancyhdr}
\usepackage{titlesec}
\usepackage{caption}
\usepackage{tabularx}
\usepackage{array}
\usepackage{float}
\usepackage{hyperref}
\usepackage[htt]{hyphenat}
\usepackage{microtype}
 
\DeclareRobustCommand{\_}{%
  \ifmmode
    \nfss@text{\textunderscore}%
  \else
    \textunderscore\allowbreak
  \fi
}
 
\renewcommand{\topfraction}{0.85}
\renewcommand{\bottomfraction}{0.85}
\renewcommand{\textfraction}{0.15}
\renewcommand{\floatpagefraction}{0.70}
 
\raggedbottom
\sloppy
 
\pagestyle{fancy}
\fancyhf{}
\fancyfoot[R]{\thepage}
\renewcommand{\headrulewidth}{0pt}
\renewcommand{\footrulewidth}{0pt}
 
\titleformat{\section}
  {\normalfont\bfseries\normalsize}
  {\thesection.}
  {0.5em}
  {}
\titleformat{\subsection}
  {\normalfont\bfseries\normalsize}
  {\thesubsection}
  {0.5em}
  {}
\titleformat{\subsubsection}
  {\normalfont\bfseries\normalsize}
  {\thesubsubsection}
  {0.5em}
  {}
 
\titlespacing*{\section}{0pt}{1.5ex plus 0.5ex minus 0.2ex}{0.8ex plus 0.2ex}
\titlespacing*{\subsection}{0pt}{1.2ex plus 0.4ex minus 0.2ex}{0.6ex plus 0.2ex}
\titlespacing*{\subsubsection}{0pt}{1.0ex plus 0.3ex minus 0.2ex}{0.5ex plus 0.2ex}
 
\DeclareCaptionLabelSeparator{hifen}{-- }
\captionsetup[figure]{
    name=Figura,
    labelsep=hifen,
    justification=centering,
    font=bf,
    labelfont=bf
}
 
\hypersetup{
    colorlinks=true,
    linkcolor=black,
    citecolor=blue,
    urlcolor=blue
}
 
\definecolor{codegreen}{rgb}{0,0.5,0}
\definecolor{codegray}{rgb}{0.45,0.45,0.45}
\definecolor{codepurple}{rgb}{0.55,0,0.75}
\definecolor{backcolour}{rgb}{0.98,0.98,0.98}
 
\lstdefinestyle{estiloc}{
    backgroundcolor=\color{backcolour},
    commentstyle=\color{codegreen}\itshape,
    keywordstyle=\color{blue}\bfseries,
    numberstyle=\tiny\color{codegray},
    stringstyle=\color{codepurple},
    basicstyle=\ttfamily\footnotesize,
    breakatwhitespace=false,
    breaklines=true,
    captionpos=b,
    keepspaces=true,
    numbers=left,
    numbersep=6pt,
    showspaces=false,
    showstringspaces=false,
    showtabs=false,
    tabsize=2,
    frame=single,
    rulecolor=\color{black!25}
}
\lstset{style=estiloc}
 
\begin{document}
 
% =========================================================================
% CAPA
% =========================================================================
\thispagestyle{empty}
\begin{center}
    {\fontsize{38}{42}\selectfont \textbf{\textsf{\textcolor[RGB]{25,25,25}{U}\hspace{1pt}\textcolor[RGB]{150,20,20}{F}\hspace{1pt}\textcolor[RGB]{200,140,20}{V}}}}\\[0.25cm]
    {\normalsize SERVIÇO PÚBLICO FEDERAL $\cdot$ MINISTÉRIO DA EDUCAÇÃO}\\[0.12cm]
    {\normalsize UNIVERSIDADE FEDERAL DE VIÇOSA $\cdot$ UFV}\\[0.12cm]
    {\normalsize CAMPUS FLORESTAL}\\[4.8cm]
 
    {\large \textbf{Trabalho 2 - PAA}}\\[0.3cm]
    {\large \textbf{Escape da Caverna com Programação Dinâmica}}\\[0.3cm]
    
 
    \begin{flushright}
        Gabriel Marcus de Oliveira Felix --- Matrícula: 5792 \\
        Otávio Soares Pedroso --- Matrícula: 6593 \\
        Pedro Mendes Gonzaga --- Matrícula: 6563 \\
        Rafaella Ferreira Pinheiro Silva --- Matrícula: 5363 \\
        Tomás Viana Resende --- Matrícula: 6556
    \end{flushright}
 
    \vfill
    {\normalsize Florestal - MG}\\[0.12cm]
    {\normalsize 2026}
\end{center}
 
\newpage
\setcounter{page}{2}
 
% =========================================================================
% Leitura e representação da caverna
% =========================================================================
 
\section{Desenvolvimento}
 
\subsection{``mapa.c'' e ``mapa.h''}
 
O módulo do mapa é responsável por abrir o arquivo de entrada, validar seu conteúdo, armazenar a caverna em memória e liberar essa memória ao final. Foi utilizado `stdio.h' para a leitura do arquivo, `stdlib.h' para a alocação dinâmica e a conversão de números, `string.h' para a comparação dos símbolos \texttt{I} e \texttt{F}, `limits.h' para validar o intervalo dos inteiros lidos e `errno.h' para detectar estouro na conversão. Os arquivos do módulo são \texttt{include/mapa.h} (interface), \texttt{src/mapa.c} (implementação) e \texttt{testes/teste\_mapa.c} (programa de teste).
 
\subsubsection*{Formato do arquivo de entrada}
 
A primeira linha do arquivo contém três inteiros separados por espaço: o número de linhas $L$, o número de colunas $C$ e os Pontos de Vida iniciais do estudante. As $L$ linhas seguintes contêm, cada uma, $C$ células separadas por espaço. Cada célula é um inteiro, que representa a variação nos Pontos de Vida ao passar por ela (negativo para monstros, positivo para poções e zero para células vazias), a letra \texttt{I} (posição inicial) ou a letra \texttt{F} (saída). O exemplo do enunciado, \texttt{Caverna1.txt}, é:
 
\begin{figure}[htbp]
    \centering
    \begin{minipage}{0.45\textwidth}
        \ttfamily\small
        \begin{tabular}{|l|}
        \hline
        4 5 40 \\
        -20 0 -20 0 0 \\
        F -10 -10 -10 -10 \\
        -20 0 0 0 +20 \\
        0 0 -10 0 I \\
        \hline
        \end{tabular}
    \end{minipage}
    \vspace{0.2cm}
    \caption{Arquivo de entrada \texttt{Caverna1.txt}}
\end{figure}
 
Esse mapa tem 4 linhas, 5 colunas e 40 Pontos de Vida iniciais, com o estudante em (3,4) e a saída em (1,0). As coordenadas começam em zero, com a linha como primeira coordenada ($x$) e a coluna como segunda ($y$), seguindo o enunciado. Valores com sinal explícito, como \texttt{+20}, são aceitos.
 
\subsubsection*{Estruturas de dados}
 
No cabeçalho \texttt{mapa.h} foi definida a estrutura principal do módulo:
 
\begin{lstlisting}[language=C, caption={Estrutura Mapa, que representa a caverna em memória.}]
typedef struct {
    int linhas, colunas;
    int hp_inicial;
    int **valor;                        /* efeito de cada celula */
    int inicio_linha, inicio_coluna;    /* posicao de I */
    int fim_linha, fim_coluna;          /* posicao de F */
} Mapa;
\end{lstlisting}
 
\begin{itemize}
    \item \texttt{linhas}, \texttt{colunas} e \texttt{hp\_inicial}: armazenam a primeira linha do arquivo;
    \item \texttt{valor}: matriz $L \times C$ alocada dinamicamente na \textit{Heap}, em que \texttt{valor[i][j]} é a variação de Pontos de Vida da célula $(i,j)$. Nas células \texttt{I} e \texttt{F} o valor armazenado é 0, pois elas não alteram os Pontos de Vida;
    \item \texttt{inicio\_linha}, \texttt{inicio\_coluna}, \texttt{fim\_linha} e \texttt{fim\_coluna}: posições de \texttt{I} e \texttt{F}, registradas durante a leitura para que não seja necessário percorrer a matriz à procura delas.
\end{itemize}
 
O cabeçalho também define os códigos de retorno da leitura, que permitem ao chamador distinguir o tipo de falha:
 
\begin{table}[htbp]
\centering
\small
\caption{Códigos de retorno da função \texttt{mapa\_ler}.}
\label{tab:codigos}
\vspace{0.2cm}
\begin{tabularx}{\textwidth}{|>{\raggedright\arraybackslash}p{4.2cm}|c|>{\raggedright\arraybackslash}X|}
\hline
\textbf{Constante} & \textbf{Valor} & \textbf{Significado} \\
\hline
\texttt{MAPA\_OK} & 0 & Leitura bem-sucedida \\
\hline
\texttt{MAPA\_ERRO\_ARQUIVO} & $-1$ & Arquivo inexistente ou ilegível \\
\hline
\texttt{MAPA\_ERRO\_MEMORIA} & $-2$ & Falha de alocação de memória \\
\hline
\texttt{MAPA\_ERRO\_FORMATO} & $-3$ & Dimensões inválidas, token inválido, células faltando, ou \texttt{I}/\texttt{F} ausente ou repetido \\
\hline
\end{tabularx}
\end{table}
 
\textbf{`mapa\_ler'}: Abre o arquivo e lê a primeira linha, validando que as dimensões são positivas. Em seguida aloca a matriz $L \times C$ e lê cada célula como texto. Se o texto for \texttt{I} ou \texttt{F}, registra a posição na estrutura e guarda valor 0; caso contrário, converte-o para inteiro com \texttt{strtol}, que aceita sinal explícito e permite detectar textos que não sejam números (como \texttt{abc}) e valores fora do intervalo de um \texttt{int}. Ao final, verifica se existe exatamente um \texttt{I} e um \texttt{F}.
 
\begin{lstlisting}[language=C, caption={Trecho central da leitura das células em mapa\_ler.}]
for (i = 0; i < mapa->linhas; i++) {
    for (j = 0; j < mapa->colunas; j++) {
        if (fscanf(f, "%63s", tok) != 1) { /* faltam celulas */
            fclose(f); mapa_liberar(mapa); return MAPA_ERRO_FORMATO;
        }
        if (strcmp(tok, "I") == 0) {
            mapa->inicio_linha = i; mapa->inicio_coluna = j;
            mapa->valor[i][j] = 0;
        } else if (strcmp(tok, "F") == 0) {
            mapa->fim_linha = i; mapa->fim_coluna = j;
            mapa->valor[i][j] = 0;
        } else if (!token_para_int(tok, &mapa->valor[i][j])) {
            fclose(f); mapa_liberar(mapa); return MAPA_ERRO_FORMATO;
        }
    }
}
\end{lstlisting}
 
\textbf{`mapa\_liberar'}: Libera cada linha da matriz, depois o vetor de ponteiros, e zera os campos da estrutura. Como os ponteiros são colocados em \texttt{NULL} ao final, a função pode ser chamada mais de uma vez sem causar erro.
 
\textbf{Decisões de implementação:}
\begin{enumerate}
    \item \textbf{Leitura dos símbolos como texto:} Como cada célula pode ser um número (inclusive com sinal, como \texttt{+20}) ou uma letra (\texttt{I} ou \texttt{F}), a leitura direta com \texttt{\%d} falharia ao encontrar uma letra. Por isso cada célula é lida como \textit{string} e só depois interpretada, o que também permite rejeitar símbolos desconhecidos em vez de ignorá-los silenciosamente.
    \item \textbf{Valor 0 em \texttt{I} e \texttt{F}:} A posição inicial e a saída não possuem monstros nem poções. Ao armazenar 0 nessas células, o valor de qualquer célula pode ser somado sem tratar esses casos separadamente.
    \item \textbf{Retorno de códigos de erro, sem \texttt{exit}:} O módulo nunca encerra o programa por conta própria. Ele devolve um código ao chamador, que decide como reagir. Isso mantém o módulo reutilizável e testável isoladamente.
    \item \textbf{Alocação linha por linha com \texttt{calloc}:} A matriz é alocada como um vetor de ponteiros e depois uma linha por vez, o que permite o acesso natural \texttt{valor[i][j]}. O uso de \texttt{calloc} deixa todos os ponteiros inicialmente nulos; assim, se a alocação falhar no meio do processo, \texttt{mapa\_liberar} consegue desfazer exatamente o que já havia sido alocado, sem vazamento de memória.
    \item \textbf{Falha segura:} Em qualquer erro, a função libera a memória já alocada e fecha o arquivo antes de retornar. Nenhuma situação de erro deixa memória pendente nem o arquivo aberto.
\end{enumerate}
 
\subsubsection*{Gerenciamento de memória}
 
A função \texttt{mapa\_ler} aloca a matriz e quem a chamou é responsável por chamar \texttt{mapa\_liberar} quando o mapa deixar de ser necessário. A hierarquia de ponteiros é representada na Figura~\ref{fig:memoria}:
 
\begin{figure}[htbp]
    \centering
    \begin{minipage}{0.7\textwidth}
        \ttfamily\small
        \begin{tabular}{|l|}
        \hline
        Mapa m (pilha) \\
        \quad linhas = 4, colunas = 5, hp\_inicial = 40 \\
        \quad valor ------> [ p0 | p1 | p2 | p3 ]   (Heap, 4 ponteiros) \\
        \quad \quad \quad \quad \quad \quad \quad  p0 --> [ -20, 0, -20, 0, 0 ] \\
        \quad \quad \quad \quad \quad \quad \quad  p1 --> [ 0, -10, -10, -10, -10 ] \\
        \quad \quad \quad \quad \quad \quad \quad  p2 --> [ -20, 0, 0, 0, 20 ] \\
        \quad \quad \quad \quad \quad \quad \quad  p3 --> [ 0, 0, -10, 0, 0 ] \\
        \hline
        \end{tabular}
    \end{minipage}
    \vspace{0.2cm}
    \caption{Alocação da matriz de valores na Heap para o exemplo \texttt{Caverna1.txt}}
    \label{fig:memoria}
\end{figure}
 
\subsubsection*{Complexidade}
 
A leitura percorre cada uma das $L \times C$ células uma única vez, realizando em cada uma operações de custo constante (leitura de um texto de tamanho limitado, comparação e conversão). Logo, a complexidade de tempo de \texttt{mapa\_ler} é $\mathcal{O}(L \times C)$. A matriz de valores ocupa $\mathcal{O}(L \times C)$ de memória, mais $\mathcal{O}(L)$ para o vetor de ponteiros, o que mantém a complexidade espacial em $\mathcal{O}(L \times C)$. A função \texttt{mapa\_liberar} percorre as $L$ linhas, com custo $\mathcal{O}(L)$.
 
% =========================================================================
% Programação Dinâmica
% =========================================================================

\subsection{``dp.c'' e ``dp.h''}

O módulo de programação dinâmica é responsável por determinar a rota que maximiza os Pontos de Vida finais do estudante ao escapar da caverna, além de identificar se existe ao menos um caminho viável que permita alcançar a saída \texttt{F} sem que a vida chegue a zero ou a valores negativos. Foram utilizadas as bibliotecas \texttt{<stdio.h>} para operações de entrada e saída, \texttt{<stdlib.h>} para alocação e liberação dinâmica de memória, e \texttt{<limits.h>} para o uso da constante \texttt{INT\_MIN}, empregada como sentinela de células inacessíveis. Os arquivos que compõem o módulo são \texttt{include/dp.h} (interface pública) e \texttt{src/dp.c} (implementação dos algoritmos).

\subsubsection*{Princípio da Programação Dinâmica e relação de recorrência}

O problema consiste em conduzir o personagem da posição inicial \texttt{I} até a saída \texttt{F}. Pelas regras estabelecidas, a partir de qualquer célula $(i,j)$ os únicos movimentos permitidos são para \textbf{cima} (linha $i-1$) e para a \textbf{esquerda} (coluna $j-1$). Dessa forma, para que o estudante alcance a célula $(i,j)$, o passo imediatamente anterior deve ter ocorrido obrigatoriamente a partir de $(i+1,j)$ (movimento para cima vindo da linha de baixo) ou de $(i,j+1)$ (movimento para a esquerda vindo da coluna à direita).

O problema exibe as duas características fundamentais que justificam o uso de Programação Dinâmica:
\begin{enumerate}
    \item \textbf{Subestrutura Ótima:} O maior valor de Pontos de Vida acumulado ao alcançar uma célula $(i,j)$ pode ser decomposto na escolha ótima entre os estados antecessores viáveis $(i+1,j)$ e $(i,j+1)$, acrescido da variação de vida da própria célula $(i,j)$;
    \item \textbf{Sobreposição de Subproblemas:} Diferentes trajetórias que partem de \texttt{I} convergem para as mesmas células intermediárias da caverna. Em vez de reavaliar recursivamente essas rotas múltiplas vezes (o que levaria a uma árvore exponencial), a tabela armazena a melhor solução para cada célula assim que é calculada.
\end{enumerate}

Sendo $\text{DP}[i][j]$ o valor máximo de Pontos de Vida com que o estudante ingressa na célula $(i,j)$, e $\text{valor}[i][j]$ a alteração de HP conferida por essa célula, a relação de recorrência que governa o preenchimento da tabela é:

\begin{equation}
\text{DP}[i][j] = \max\Big(\text{DP}[i+1][j],\; \text{DP}[i][j+1]\Big) + \text{valor}[i][j]
\label{eq:recorrencia_dp}
\end{equation}

sob as seguintes condições de contorno e restrições de viabilidade:
\begin{itemize}
    \item \textbf{Caso Base:} Na posição inicial $I = (i_I, j_I)$, o valor é inicializado diretamente com o HP fornecido pelo arquivo de entrada:
    $$\text{DP}[i_I][j_I] = \text{hp\_inicial}$$
    \item \textbf{Condição de Vida Positiva para Transição:} Um movimento a partir de um vizinho só é considerado se o personagem estiver vivo naquela célula, ou seja, $\text{DP}[i+1][j] > 0$ ou $\text{DP}[i][j+1] > 0$. Vizinhos não visitados ou onde houve morte contêm \texttt{INT\_MIN} e são desconsiderados;
    \item \textbf{Sobrevivência na Célula Atual:} Caso o valor resultante da transição $\text{DP}[i][j] \le 0$, o estudante perece naquela sala, tornando-a inviável como ponto de passagem. Nesses casos, o valor não é consolidado como rota válida e o estado permanece como inacessível.
\end{itemize}

\subsubsection*{Ordem topológica de preenchimento}

Como a transição da célula $(i,j)$ depende dos valores já computados em $(i+1,j)$ e $(i,j+1)$, o algoritmo deve resolver primeiro as linhas inferiores e colunas à direita. O algoritmo inicia na linha de \texttt{I} e percorre a matriz em sentido decrescente:

\begin{itemize}
    \item O laço externo itera a linha $i$ de $i_I$ até $0$ (de baixo para cima);
    \item O laço interno itera a coluna $j$ de $j_I$ até $0$ (da direita para a esquerda).
\end{itemize}

Dessa forma, no momento em que qualquer célula $(i,j)$ é avaliada, suas dependências $(i+1,j)$ e $(i,j+1)$ já foram completamente resolvidas e memorizadas na tabela.

\subsubsection*{Estruturas de dados}

O cabeçalho \texttt{include/dp.h} define a estrutura central de resolução:

\begin{lstlisting}[language=C, caption={Estrutura Solucao que consolida os resultados da Programação Dinâmica.}]
struct Solucao {
    int existe_caminho;       /* 1 se F e alcancavel com HP > 0, 0 caso contrario */
    int possivel;             /* mantido para retrocompatibilidade */
    int hp_final;             /* maior HP ao chegar em F */
    int **melhor_hp;          /* tabela L x C com o maior HP em cada celula */
    Coordenada **pai;         /* tabela L x C com a celula antecessora */
};
\end{lstlisting}

\begin{itemize}
    \item \texttt{existe\_caminho} e \texttt{possivel}: \textit{flags} booleanas que registram se a saída \texttt{F} foi alcançada com vida positiva ($\text{HP} > 0$);
    \item \texttt{hp\_final}: armazena o valor ótimo de Pontos de Vida acumulados na célula de saída $(fim\_linha, fim\_coluna)$;
    \item \texttt{melhor\_hp}: matriz $L \times C$ de inteiros alocada na \textit{Heap}, correspondente à tabela de programação dinâmica, onde cada entrada guarda o maior HP viável;
    \item \texttt{pai}: matriz $L \times C$ de coordenadas (\texttt{Coordenada \{int linha; int coluna;\}}), onde cada posição guarda a célula imediatamente anterior que originou o caminho ótimo até $(i,j)$. Essa matriz viabiliza a posterior reconstrução do trajeto sem necessidade de recalcular a DP.
\end{itemize}

\subsubsection*{Funções e implementação}

O módulo organiza o fluxo de resolução em duas funções públicas e quatro rotinas estáticas auxiliares de gerenciamento de matrizes:

\textbf{`alocar\_matriz\_int' e `alocar\_matriz\_coord'}: Alocam matrizes bidimensionais na \textit{Heap} utilizando \texttt{calloc}, garantindo inicialização limpa. Caso ocorra falta de memória no meio do laço de alocação das linhas, a função desaloca imediatamente as linhas previamente concedidas antes de retornar \texttt{NULL}, evitando qualquer vazamento residual de memória.

\textbf{`liberar\_matriz\_int' e `liberar\_matriz\_coord'}: Percorrem as $L$ linhas da matriz liberando cada bloco e, ao término, liberam o vetor principal de ponteiros.

\textbf{`dp\_resolver'}: Orquestra a execução da programação dinâmica:
\begin{enumerate}
    \item Inicializa os campos da estrutura \texttt{Solucao} com valores nulos e aloca as matrizes \texttt{melhor\_hp} e \texttt{pai};
    \item Preenche todas as células de \texttt{melhor\_hp} com a constante sentinela \texttt{INT\_MIN} e as coordenadas em \texttt{pai} com $(-1,-1)$;
    \item Define o caso base na célula inicial $(inicio\_linha, inicio\_coluna)$ com o HP inicial lido do mapa;
    \item Executa os laços aninhados de programação dinâmica em ordem decrescente de linhas e colunas;
    \item Para cada célula, examina se o vizinho de baixo $(i+1,j)$ e o da direita $(i,j+1)$ possuem HP positivo, elegendo aquele com maior valor acumulado;
    \item Soma o efeito da célula $\text{valor}[i][j]$. Se o resultado for estritamente positivo ($\text{HP} > 0$), atualiza a tabela e define o predecessor;
    \item Ao final do preenchimento, inspeciona a célula de destino $(fim\_linha, fim\_coluna)$. Se ela foi alcançada com $\text{HP} > 0$, marca a solução como viável, atribui o HP final e retorna $1$. Caso contrário, define que não existe caminho e retorna $0$.
\end{enumerate}

\begin{lstlisting}[language=C, caption={Trecho central do preenchimento da tabela em dp\_resolver.}]
for (i = il; i >= 0; i--) {
    for (j = ic; j >= 0; j--) {
        if (i == il && j == ic)
            continue; /* caso base ja tratado */

        int melhor = INT_MIN;
        Coordenada origem = {-1, -1};

        /* Veio de baixo (movimento para cima) */
        if (i + 1 < mapa->linhas && solucao->melhor_hp[i + 1][j] > 0) {
            if (solucao->melhor_hp[i + 1][j] > melhor) {
                melhor = solucao->melhor_hp[i + 1][j];
                origem.linha = i + 1;
                origem.coluna = j;
            }
        }

        /* Veio da direita (movimento para esquerda) */
        if (j + 1 < mapa->colunas && solucao->melhor_hp[i][j + 1] > 0) {
            if (solucao->melhor_hp[i][j + 1] > melhor) {
                melhor = solucao->melhor_hp[i][j + 1];
                origem.linha = i;
                origem.coluna = j + 1;
            }
        }

        if (melhor != INT_MIN) {
            int hp_novo = melhor + mapa->valor[i][j];
            if (hp_novo > 0) {
                solucao->melhor_hp[i][j] = hp_novo;
                solucao->pai[i][j] = origem;
            }
        }
    }
}
\end{lstlisting}

\textbf{`dp\_liberar'}: Desaloca com segurança as duas matrizes dinâmicas da solução chamando as rotinas de liberação, zera todos os ponteiros para \texttt{NULL} e reseta os contadores da estrutura, permitindo que a função seja invocada repetidas vezes sem risco de falha de segmentação (\textit{double free}).

\textbf{Decisões de implementação:}
\begin{enumerate}
    \item \textbf{Diferenciação entre célula inalcançável e morte:} Ao utilizar \texttt{INT\_MIN} como valor sentinela, distingue-se com exatidão uma célula que nunca pôde ser alcançada por nenhuma rota daquelas em que o personagem alcança mas tem vida nula ou negativa ($\le 0$).
    \item \textbf{Prevenção contra transições a partir de personagens mortos:} A condicional \texttt{melhor\_hp[...] > 0} garante que caminhos onde o estudante morre não continuem a ser propagados na tabela, respeitando rigorosamente a restrição de que a vida deve ser sempre positiva.
    \item \textbf{Gravação implícita de predecessores para reconstrução:} O registro da coordenada de origem em \texttt{pai[i][j]} durante o próprio preenchimento descarta a necessidade de manter pilhas de busca ou executar nova exploração retroativa comparando diferenças de valores, reduzindo a complexidade de recuperação do percurso para tempo linear estrito.
    \item \textbf{Robustez e tratamento defensivo:} Caso ponteiros nulos sejam fornecidos ou as coordenadas de início/fim estejam fora dos limites do mapa, a rotina aborta a resolução de forma controlada devolvendo código de erro sem interrupções abruptas (\textit{crash}).
\end{enumerate}

\subsubsection*{Gerenciamento de memória}

A estrutura \texttt{Solucao} reside normalmente na pilha da função chamadora, enquanto seus membros \texttt{melhor\_hp} e \texttt{pai} apontam para blocos alocados dinamicamente na \textit{Heap}. A Figura~\ref{fig:memoria_dp} ilustra a disposição dos ponteiros e matrizes para o exemplo do enunciado:

\begin{figure}[htbp]
    \centering
    \begin{minipage}{0.78\textwidth}
        \ttfamily\small
        \begin{tabular}{|l|}
        \hline
        Solucao s (pilha) \\
        \quad existe\_caminho = 1, hp\_final = 50 \\
        \quad melhor\_hp -----> [ p0 | p1 | p2 | p3 ] (Heap, 4 ponteiros de int*) \\
        \quad \quad \quad \quad \quad \quad \quad p0 --> [ --, --, --, --, -- ]  (-- representa INT\_MIN) \\
        \quad \quad \quad \quad \quad \quad \quad p1 --> [ 50, 50, --, --, -- ]  (F em (1,0) = 50) \\
        \quad \quad \quad \quad \quad \quad \quad p2 --> [ --, 60, 60, 60, 60 ] \\
        \quad \quad \quad \quad \quad \quad \quad p3 --> [ --, --, --, --, 40 ]  (I em (3,4) = 40) \\[0.15cm]
        \quad pai -----------> [ q0 | q1 | q2 | q3 ] (Heap, 4 ponteiros de Coord*) \\
        \quad \quad \quad \quad \quad \quad \quad q0 --> [ (-1,-1), (-1,-1), \dots ] \\
        \quad \quad \quad \quad \quad \quad \quad q1 --> [ (1,1), (2,1), \dots ]  (origem de F é (1,1)) \\
        \quad \quad \quad \quad \quad \quad \quad q2 --> [ (-1,-1), (2,2), (2,3), (2,4), (3,4) ] \\
        \quad \quad \quad \quad \quad \quad \quad q3 --> [ \dots, (-1,-1) ]       (I não tem predecessor) \\
        \hline
        \end{tabular}
    \end{minipage}
    \vspace{0.2cm}
    \caption{Alocação e valores das tabelas \texttt{melhor\_hp} e \texttt{pai} na Heap para \texttt{Caverna1.txt}}
    \label{fig:memoria_dp}
\end{figure}

\subsubsection*{Complexidade}

A função \texttt{dp\_resolver} visita cada uma das $L \times C$ células no máximo uma vez. Em cada célula, são realizadas no máximo duas consultas a vizinhos adjacentes, comparações de inteiros e uma soma aritmética, todas operações de custo elementar $\mathcal{O}(1)$. Portanto, a complexidade de tempo de resolução por Programação Dinâmica é rigorosamente:
$$\mathcal{O}(L \times C)$$
Em termos de espaço, são alocadas duas matrizes de dimensão $L \times C$ (\texttt{melhor\_hp} de inteiros e \texttt{pai} de estruturas \texttt{Coordenada}), além dos vetores auxiliares de ponteiros de dimensão $L$, totalizando complexidade espacial de:
$$\mathcal{O}(L \times C)$$
A liberação de memória em \texttt{dp\_liberar} percorre as $L$ linhas das duas matrizes, demandando tempo $\mathcal{O}(L)$ e espaço adicional $\mathcal{O}(1)$.



\section{Resultados}
 
\subsection{Testes do módulo de leitura}
 
Para validar o módulo de forma isolada, foi criado o programa \texttt{teste\_mapa}, que lê o arquivo informado, imprime as dimensões, os Pontos de Vida iniciais, as posições de \texttt{I} e \texttt{F} e a matriz lida, e então libera a memória. Ele é compilado com o comando \texttt{make} dentro da pasta \texttt{tp2} (que executa \texttt{gcc -Wall -Wextra -std=c99 -Iinclude -o teste\_mapa testes/teste\_mapa.c src/mapa.c}, sem advertências) e executado com o caminho de um arquivo de caverna. A execução com o exemplo do enunciado é mostrada na Figura~\ref{fig:teste_ok}:
 
\begin{figure}[htbp]
    \centering
    \begin{minipage}{0.8\textwidth}
        \ttfamily\small
        \begin{tabular}{|l|}
        \hline
        \$ ./teste\_mapa testes/Caverna1.txt \\
        Dimensoes: 4 x 5 | HP inicial: 40 \\
        I = (3, 4) | F = (1, 0) \\
        \ -20 \ \ 0 \ -20 \ \ 0 \ \ 0 \\
        \ \ \ 0 \ -10 \ -10 \ -10 \ -10 \\
        \ -20 \ \ 0 \ \ 0 \ \ 0 \ \ 20 \\
        \ \ \ 0 \ \ 0 \ -10 \ \ 0 \ \ 0 \\
        \hline
        \end{tabular}
    \end{minipage}
    \vspace{0.2cm}
    \caption{Leitura correta do arquivo \texttt{Caverna1.txt}}
    \label{fig:teste_ok}
\end{figure}
 
A saída confere com o arquivo de entrada: as dimensões e os Pontos de Vida iniciais estão corretos, o \texttt{I} foi localizado em $(3,4)$ e o \texttt{F} em $(1,0)$, os valores de todas as células foram preservados, o valor \texttt{+20} foi interpretado como 20, e as células \texttt{I} e \texttt{F} foram armazenadas com valor 0.
 
Foram testadas também entradas inválidas, verificando se o módulo devolve o código de erro apropriado sem encerrar o programa. As execuções estão na Figura~\ref{fig:teste_erro}, e a Tabela~\ref{tab:casos} resume todos os casos:
 
\begin{figure}[htbp]
    \centering
    \begin{minipage}{0.8\textwidth}
        \ttfamily\small
        \begin{tabular}{|l|}
        \hline
        \$ ./teste\_mapa nao\_existe.txt \\
        mapa\_ler falhou (codigo -1) \\[0.15cm]
        \$ cat testes/sem\_I.txt \\
        2 2 10 \\
        F 1 \\
        1 1 \\
        \$ ./teste\_mapa testes/sem\_I.txt \\
        mapa\_ler falhou (codigo -3) \\
        \hline
        \end{tabular}
    \end{minipage}
    \vspace{0.2cm}
    \caption{Tratamento de arquivo inexistente e de mapa sem a posição inicial}
    \label{fig:teste_erro}
\end{figure}
 
\begin{table}[htbp]
\centering
\small
\caption{Casos de teste do módulo de leitura.}
\label{tab:casos}
\vspace{0.2cm}
\begin{tabularx}{\textwidth}{|>{\raggedright\arraybackslash}p{3.3cm}|>{\raggedright\arraybackslash}X|>{\centering\arraybackslash}p{2.4cm}|>{\centering\arraybackslash}p{2.0cm}|}
\hline
\textbf{Caso} & \textbf{Descrição} & \textbf{Esperado} & \textbf{Obtido} \\
\hline
\texttt{Caverna1.txt} & Mapa $4 \times 5$ do enunciado, HP 40, com \texttt{+20} & Leitura correta (código 0) & Conforme \\
\hline
\texttt{nao\_existe.txt} & Arquivo que não existe & Código $-1$ & Conforme \\
\hline
\texttt{sem\_I.txt} & Mapa $2 \times 2$ sem a posição inicial \texttt{I} & Código $-3$ & Conforme \\
\hline
\end{tabularx}
\end{table}
 
Os arquivos de entrada utilizados estão na pasta \texttt{tp2/testes} do repositório: \texttt{Caverna1.txt} e \texttt{sem\_I.txt}.
 
\end{document}